\section{String Theory from Observer Patch Holography}\label{string-theory-from-observer-patch-holography}

\begin{quote}
OPH is an observer-centric framework for fundamental physics. From one calibrated constant (the pixel area), structural axioms (A1-A4 + MAR), and one cosmological input (screen capacity), OPH derives scaling-limit semiclassical gravity, compact gauge reconstruction, particle mass hierarchy, and key quantitative outputs.
\end{quote}

\subsection{A Complete Derivation}\label{a-complete-derivation}

This document provides a self-contained derivation showing that Observer Patch Holography (OPH) necessarily implies a string-theoretic description of its fundamental degrees of freedom. The derivation uses only OPH axioms and standard mathematical results,no external string theory inputs are required.

\begin{center}\rule{0.5\linewidth}{0.5pt}\end{center}

\section{1. Overview: What We Derive}\label{1-overview-what-we-derive}

We establish the following chain of implications using only OPH axioms:

\begin{verbatim}
OPH Axioms (A1-A4) + Regulator Premises (R0-R1)
    ↓
Edge-Center Decomposition + Heat-Kernel Sector Weights
    ↓
2D Yang-Mills Theory on the Holographic Screen
    ↓
String Theory (Worldsheet Expansion via Gross-Taylor)
\end{verbatim}

Additionally, we show that OPH supplies:

\begin{itemize}
\tightlist
\item
  \textbf{Worldsheet kinematics}: Conformal/modular sewing from geometric modular flow
\item
  \textbf{Target-space dynamics}: conditional scaling-limit Einstein branch from entanglement equilibrium
\item
  \textbf{Higher gauge structure}: The natural slot for B-field/gerbe data
\end{itemize}

This establishes that OPH contains string theory as an emergent description of its edge/boundary degrees of freedom.

\begin{center}\rule{0.5\linewidth}{0.5pt}\end{center}

\section{2. OPH Inputs Used}\label{2-oph-inputs-used}

We use exactly the following OPH ingredients from the main paper and technical supplement:

\subsection{2.1 Core Axioms}\label{21-core-axioms}

\textbf{Axiom A1 (Screen Net):} Physical reality is encoded on a horizon screen \(S^2\) carrying a net of von Neumann algebras \(P \mapsto \mathcal{A}(P)\) for patches \(P \subset S^2\).

\textbf{Axiom A2 (Overlap Consistency):} For overlapping patches \(P_1 \cap P_2 \neq \emptyset\), the restrictions of local states must agree: \[\omega_1|_{\mathcal{A}(P_1 \cap P_2)} = \omega_2|_{\mathcal{A}(P_1 \cap P_2)}\]

\textbf{Axiom A3 (Generalized Entropy):} A generalized entropy functional exists: \[S_{\text{gen}}(C) = \frac{A(\partial C)}{4G} + S_{\text{bulk}}(C)\] satisfying quantum focusing (monotonicity under inclusion along null generators).

\textbf{Axiom A4 (Local Markov/Recoverability):} For tripartitions \(A\)-\(B\)-\(D\) across separators, the conditional mutual information is small: \[I(A:D|B) \leq \varepsilon\] and recovery maps exist with controlled error.

\subsection{2.2 Regulator Premises}\label{22-regulator-premises}

\textbf{Premise R0 (Type-I Local Algebras):} At the UV regulator scale, local patch algebras are type-I factors (finite-dimensional matrix algebras).

\textbf{Premise R1 (Boundary Gauge Invariance):} For a region \(R \subset S^2\), the physical algebra is: \[\mathcal{A}(R) = \mathcal{B}(\tilde{\mathcal{H}}_R)^{G_{\partial R}}\] where \(G_{\partial R}\) is the boundary gauge group.

\subsection{2.3 Key Derived Results}\label{23-key-derived-results}

\textbf{Edge-Center Decomposition (Theorem 1.1):} Under R0 and R1, the collar Hilbert space decomposes as: \[\mathcal{H}_{B_\delta} = \bigoplus_{\alpha} \left(\mathcal{H}_{b_L^\alpha} \otimes \mathcal{H}_{b_R^\alpha}\right)\]

\textbf{Heat-Kernel Sector Weights (Theorem 5.1):} Under MaxEnt with bi-invariant constraints: \[p_R(t) \propto d_R \, e^{-t \, C_2(R)}\] where \(d_R = \dim R\) and \(C_2(R)\) is the quadratic Casimir.

\textbf{Geometric Modular Flow (Theorem 2.2):} Cap modular flow acts as the unique cap-preserving conformal dilation with period \(2\pi\).

\begin{center}\rule{0.5\linewidth}{0.5pt}\end{center}

\section{3. Einstein Equation from Entanglement Equilibrium}\label{3-einstein-equation-from-entanglement-equilibrium}

Before deriving string theory, we first establish that OPH recovers gravity,the "target-space dynamics" that any string theory must reproduce.

\subsection{3.1 Setup: Null Deformation Problem}\label{31-setup-null-deformation-problem}

Pick a local causal diamond/cap region \(C\) in the emergent bulk. Focus on one null sheet of \(\partial C\), generated by a future-directed null vector field \(k^a\) with affine parameter \(\lambda\).

Consider an infinitesimal deformation moving the cut along the null generators: \[\Sigma \to \Sigma(\lambda)\]

\subsection{3.2 Generalized Entropy Variation}\label{32-generalized-entropy-variation}

By A3, generalized entropy is: \[S_{\text{gen}}(\Sigma) = \frac{A(\Sigma)}{4G} + S_{\text{bulk}}(\Sigma)\]

Its first variation under the null deformation: \[\delta S_{\text{gen}} = \delta\left(\frac{A}{4G}\right) + \delta S_{\text{bulk}}\]

\textbf{Entanglement equilibrium premise:} At the MaxEnt reference state, \(\delta S_{\text{gen}} = 0\) for allowed local variations.

\subsection{3.3 Area Variation via Raychaudhuri}\label{33-area-variation-via-raychaudhuri}

Let \(\theta(\lambda)\) be the expansion of the null congruence. The Raychaudhuri equation: \[\frac{d\theta}{d\lambda} = -\frac{1}{2}\theta^2 - \sigma_{ab}\sigma^{ab} - R_{ab}k^a k^b\]

At equilibrium (\(\theta(0) = 0\), \(\sigma(0) = 0\)): \[\left.\frac{d\theta}{d\lambda}\right|_{\lambda=0} = -R_{kk}\]

Therefore: \[\delta\left(\frac{A}{4G}\right) = -\frac{1}{4G}\int d^2x \sqrt{h_0} \int d\lambda \, \lambda \, R_{kk} + \cdots\]

\subsection{3.4 Bulk Entropy Variation via Modular Theory}\label{34-bulk-entropy-variation-via-modular-theory}

\textbf{Entanglement First Law:} For small perturbations around reference state \(\omega\): \[\delta S_{\text{bulk}} = \delta\langle K \rangle\] where \(K = -\log \rho_C\) is the modular Hamiltonian.

\textbf{Geometric Modular Hamiltonian (N1-N3):} For null-deformed regions: \[K = 2\pi \int d^2x \int d\lambda \, \lambda \, T_{kk}(\lambda, x) + K_\partial + O(\varepsilon)\]

Therefore: \[\delta S_{\text{bulk}} = 2\pi \int d^2x \int d\lambda \, \lambda \, \delta\langle T_{kk} \rangle + \cdots\]

\subsection{3.5 Entanglement Equilibrium Implies Einstein}\label{35-entanglement-equilibrium-implies-einstein}

Setting \(\delta S_{\text{gen}} = 0\): \[0 = -\frac{1}{4G}\int d^2x \int d\lambda \, \lambda \, R_{kk} + 2\pi \int d^2x \int d\lambda \, \lambda \, \langle T_{kk} \rangle\]

Since this holds for arbitrary localized deformations: \[R_{kk} = 8\pi G \langle T_{kk} \rangle \quad \text{for all null } k^a\]

\subsection{3.6 Upgrade to Full Einstein Equation}\label{36-upgrade-to-full-einstein-equation}

Define \(E_{ab} := R_{ab} - 8\pi G \langle T_{ab} \rangle\).

The result \(R_{kk} = 8\pi G T_{kk}\) for all null \(k\) means: \[E_{ab} k^a k^b = 0 \quad \forall \text{ null } k\]

\textbf{Lemma (Null Data Ambiguity):} If \(X_{ab} k^a k^b = 0\) for all null \(k\), then \(X_{ab} = \phi \, g_{ab}\).

Therefore \(E_{ab} = \phi \, g_{ab}\), giving: \[R_{ab} = 8\pi G \langle T_{ab} \rangle + \phi \, g_{ab}\]

In Einstein tensor form: \[\boxed{G_{ab} + \Lambda g_{ab} = 8\pi G \langle T_{ab} \rangle}\]

where \(\Lambda := \frac{1}{2}R - \phi\) is constant by the Bianchi identity.

\textbf{Key point:} \(\Lambda\) is not fixed by local null modular data,it requires global input (screen capacity).

\begin{center}\rule{0.5\linewidth}{0.5pt}\end{center}

\section{4. Edge Sectors and 2D Yang-Mills}\label{4-edge-sectors-and-2d-yang-mills}

\subsection{4.1 Edge-Sector Structure}\label{41-edge-sector-structure}

From the OPH edge-center decomposition, the collar Hilbert space has structure: \[\mathcal{H}_{\text{collar}} \cong \bigoplus_{\alpha} \left(\mathcal{H}_{b_L^\alpha} \otimes \mathcal{H}_{b_R^\alpha}\right)\]

with \(\alpha\) labeling irreps/charges of the boundary gauge group.

This is exactly the structure of gauge theory with boundary: cutting gauge space produces edge modes whose labels are representation data.

\subsection{4.2 Heat-Kernel Weights from MaxEnt}\label{42-heat-kernel-weights-from-maxent}

OPH\textquotesingle s MaxEnt + bi-invariant constraints give sector probabilities: \[p_R(t) \propto d_R \, e^{-t \, C_2(R)}\]

Because of the Peter-Weyl doubling (loops see both factors), the effective weight is: \[w_R(t) \propto d_R \, p_R(t) \propto d_R^2 \, e^{-t \, C_2(R)}\]

\subsection{4.3 Edge Partition Function}\label{43-edge-partition-function}

Define the edge partition function at diffusion parameter \(t\): \[Z_{\text{edge}}(t) := \sum_R d_R^2 \, e^{-t \, C_2(R)}\]

\subsection{4.4 Identification with Group Heat Kernel}\label{44-identification-with-group-heat-kernel}

The heat kernel on compact group \(G\), expanded in characters: \[K_t(U) = \sum_R d_R \, \chi_R(U) \, e^{-t C_2(R)}\]

At the identity \(U = 1\) (using \(\chi_R(1) = d_R\)): \[K_t(1) = \sum_R d_R^2 \, e^{-t C_2(R)} = Z_{\text{edge}}(t)\]

\textbf{Result:} The OPH edge ensemble is literally the heat kernel at the identity.

\subsection{4.5 Gluing = Markov = Area Additivity}\label{45-gluing--markov--area-additivity}

Heat kernels satisfy the Chapman-Kolmogorov (gluing) property: \[K_{t_1 + t_2}(U) = \int_G dV \, K_{t_1}(UV^{-1}) \, K_{t_2}(V)\]

This matches OPH patch sewing: add diffusion parameters when gluing collars/strips.

\subsection{4.6 Exact Match with 2D Yang-Mills}\label{46-exact-match-with-2d-yang-mills}

The partition function of 2D Yang-Mills on a closed genus-\(g\) surface \(\Sigma_g\) with area \(A\): \[Z_{\text{2D YM}}(\Sigma_g) = \sum_R (d_R)^{2-2g} \exp\left[-\frac{g_{\text{YM}}^2 A}{2} C_2(R)\right]\]

For \(\Sigma_0 = S^2\) (genus 0): \[Z_{\text{2D YM}}(S^2) = \sum_R d_R^2 \exp\left[-\frac{g_{\text{YM}}^2 A}{2} C_2(R)\right]\]

\textbf{Identification:} \[t = \frac{g_{\text{YM}}^2 A}{2}\]

\textbf{Theorem 4.1:} OPH edge-sector MaxEnt gives a fully-fledged 2D Yang-Mills theory on the screen, with collar diffusion parameter \(t\) playing the role of "coupling × area".

\begin{center}\rule{0.5\linewidth}{0.5pt}\end{center}

\section{5. 2D Yang-Mills as String Theory}\label{5-2d-yang-mills-as-string-theory}

\subsection{5.1 The Gross-Taylor Expansion}\label{51-the-gross-taylor-expansion}

In the large-\(N\) limit for \(G = SU(N)\), 2D Yang-Mills admits an exact reorganization as a sum over worldsheets (Gross-Taylor, 1993).

The partition function expands as: \[\log Z = \sum_{h \geq 0} N^{2-2h} F_h(\lambda), \quad \lambda := g_{\text{YM}}^2 N\]

where:

\begin{itemize}
\tightlist
\item
  \(h\) is the worldsheet genus
\item
  \(g_s \sim 1/N\) is the string coupling
\item
  Contributions scale as \(N^{2-2h}\) for worldsheet genus \(h\)
\end{itemize}

\subsection{5.2 Worldsheet Interpretation}\label{52-worldsheet-interpretation}

The \(1/N\) expansion coefficients count \textbf{branched coverings} of the target surface,i.e., "string configurations without folds."

This is a genuine string theory:

\begin{itemize}
\tightlist
\item
  \textbf{Target space:} The 2D screen surface
\item
  \textbf{Worldsheet:} The covering surface summed over
\item
  \textbf{String coupling:} \(g_s \sim 1/N\)
\end{itemize}

\subsection{\texorpdfstring{5.3 The Large-\(N\) Regime}{5.3 The Large-N Regime}}\label{53-the-large-n-regime}

The Gross-Taylor expansion is controlled in the large-\(N\) limit. OPH provides an effective large-\(N\) parameter from screen microstructure: many independent edge channels/pixels whose combined algebra behaves like large matrix degrees of freedom. The effective edge Hilbert space dimension scales as: \[\dim(\tilde{\mathcal{H}}_\partial) \sim N^{\#(\text{boundary cells})}\]

so \(1/N\) becomes the handle-counting coupling.

\subsection{5.4 Genus Expansion}\label{54-genus-expansion}

Define generating functional: \[F(t) := \log Z_{\text{edge}}(t)\]

In controlled large-\(N\) regime: \[F(t) = \sum_{h \geq 0} N^{2-2h} F_h(\lambda), \quad \lambda := g_{\text{YM}}^2 N\]

This is a perturbative closed-string genus expansion.

\begin{center}\rule{0.5\linewidth}{0.5pt}\end{center}

\section{6. Worldsheet Kinematics from Modular Flow}\label{6-worldsheet-kinematics-from-modular-flow}

\subsection{6.1 Conformal Structure}\label{61-conformal-structure}

OPH provides conformal structure from two sources:

\textbf{Global conformal symmetry:} \[\text{Conf}^+(S^2) \cong \text{PSL}(2,\mathbb{C}) \cong \text{SO}^+(3,1)\]

This gives Lorentz kinematics,interpreting screen angles as null directions of 3+1D spacetime.

\textbf{Geometric modular flow:} Cap modular flow is forced to be conformal dilation with period \(2\pi\) (Euclidean regularity).

\subsection{6.2 Sewing Structure}\label{62-sewing-structure}

String worldsheets are sewn along boundaries/cuts. Consistent sewing requires conformal/modular data.

OPH\textquotesingle s Markov/additivity theorems on collars derive this exact factorization structure that makes sewing well-behaved.

\subsection{6.3 The Analogy with Standard String Theory}\label{63-the-analogy-with-standard-string-theory}

{\def\LTcaptype{none} % do not increment counter
\begin{longtable}[]{@{}>{\RaggedRight\arraybackslash}p{0.480\textwidth}>{\RaggedRight\arraybackslash}p{0.480\textwidth}@{}}
\toprule\noalign{}
String Theory & OPH \\
\midrule\noalign{}
\endhead
\bottomrule\noalign{}
\endlastfoot
Worldsheet conformal symmetry & Screen conformal group \\
Modular invariance & Overlap consistency \\
Target-space Einstein eq. & Entanglement equilibrium \\
Sewing axioms & Markov collar gluing \\
\end{longtable}
}

\begin{center}\rule{0.5\linewidth}{0.5pt}\end{center}

\section{7. The B-Field Slot: Higher Gauge Structure}\label{7-the-b-field-slot-higher-gauge-structure}

\subsection{7.1 Nonabelian Gluing Cocycles}\label{71-nonabelian-gluing-cocycles}

OPH includes gluing by a crossed-module 2-cocycle \((g_{ij}, h_{ijk})\) on triple overlaps, with coherence on quadruple overlaps.

This is precisely the higher-gauge/gerbe structure associated in string theory with:

\begin{itemize}
\tightlist
\item
  2-form gauge field \(B\) (Kalb-Ramond field)
\item
  Gauge invariance \(B \to B + d\eta\)
\end{itemize}

\subsection{7.2 Strictification and Trivialization}\label{72-strictification-and-trivialization}

OPH shows strictification exists iff the abelian \v{C}ech \(H^2\) class is trivial.

This parallels: "background \(B\)-field is a gerbe class; strings couple to it; trivialization conditions matter."

\subsection{7.3 D-Brane Interpretation}\label{73-d-brane-interpretation}

Certain defects/boundaries where the class trivializes are the seed of D-brane-like boundary conditions.

The mathematical structure where string theory places B-field/gerbe data is directly present in OPH. The full open/closed string boundary state formalism follows from extending this framework to include D-brane-like boundary conditions.

\begin{center}\rule{0.5\linewidth}{0.5pt}\end{center}

\section{8. OPH String Scale and Parameters}\label{8-oph-string-scale-and-parameters}

\subsection{8.1 String Scale from Pixel Area}\label{81-string-scale-from-pixel-area}

OPH has a UV length from pixel area: \[a_{\text{cell}} \approx 1.63 \, \ell_P^2, \quad \ell_{\text{UV}} = \sqrt{a_{\text{cell}}} \approx 1.28 \, \ell_P\]

Natural identification: \[\alpha' \sim \ell_{\text{UV}}^2 \sim a_{\text{cell}}\]

String tension: \[T_{\text{string}} = \frac{1}{2\pi\alpha'} \sim \frac{1}{2\pi a_{\text{cell}}}\]

Numerically (in Planck units): \[T_{\text{string}} \approx \frac{1}{2\pi \cdot 1.63} \, \ell_P^{-2} \approx 0.098 \, \ell_P^{-2}\]

This is order-Planck tension, as expected if UV cutoff sits near \(\ell_P\).

\subsection{8.2 Worldsheet YM Coupling}\label{82-worldsheet-ym-coupling}

From OPH: \[g_{\text{ent}}^2 = \frac{t}{2\pi}\]

In 2D YM, heat-kernel time is \(t \sim (g_{\text{YM}}^2 A)/2\).

So OPH\textquotesingle s \(t\) really does play the role of 2D YM area-time / diffusion parameter.

\subsection{8.3 Current Algebra Central Charge (WZW)}\label{83-current-algebra-central-charge-wzw}

If edge gauge group embeds into worldsheet current algebra (WZW-type): \[c = \frac{k \, \dim\mathfrak{g}}{k + h^\vee}\]

For \(SU(3) \times SU(2) \times U(1)\) with \(h^\vee_{SU(3)} = 3\), \(h^\vee_{SU(2)} = 2\): \[c_{\text{int}}(k) = \frac{8k}{k+3} + \frac{3k}{k+2} + 1\]

For 4D superstring compactification with internal \(c_{\text{int}} \approx 9\): \[k = \frac{5 + \sqrt{89}}{2} \approx 7.22\]

Integer levels \(k = 7\) or \(8\) bracket \(c_{\text{int}} \approx 9\).

\textbf{Conclusion:} OPH gauge content is numerically consistent with reasonable-level current algebra sector inside critical worldsheet theory.

\begin{center}\rule{0.5\linewidth}{0.5pt}\end{center}

\section{9. From 2D Yang-Mills to Critical Superstrings}\label{9-from-2d-yang-mills-to-critical-superstrings}

The derivation chain from OPH to string theory proceeds in two stages.

\subsection{9.1 Established Results}\label{91-established-results}

The following are derived from the OPH axioms (A1--A4, R0--R1, MAR):

\begin{enumerate}
\def\labelenumi{\arabic{enumi}.}
\tightlist
\item
  Finite microscopic DoF on the screen, type-I local algebras, boundary gauge invariance
\item
  EC decomposition and Markov structure on collars
\item
  Sector category and reconstruction of the compact gauge group, with \(SU(3) \times SU(2) \times U(1)/\mathbb{Z}_6\) selected under the extended MAR admissibility package (see Part II of this manuscript)
\item
  Heat-kernel edge-sector weights \(p_R \propto d_R e^{-t C_2(R)}\)
\item
  Identification with the 2D Yang-Mills partition function (Theorem 4.1)
\item
  Geometric modular Hamiltonians with \(2\pi\) Euclidean regularity and Lorentz symmetry from Conf\((S^2)\) in the refinement/scaling regime
\item
  Conditional scaling-limit Einstein branch from entanglement equilibrium
\item
  Nonabelian gluing (crossed module/2-group) providing higher gauge data
\end{enumerate}

The genus expansion \(\log Z = \sum_{h \geq 0} N^{2-2h} F_h(\lambda)\) follows from the Gross-Taylor rewriting of 2D Yang-Mills in the large-\(N\) regime, with \(N\) identified as the effective edge Hilbert space parameter.

\subsection{9.2 Extensions to Critical Superstrings}\label{92-extensions-to-critical-superstrings}

The full path from OPH to critical superstring theory requires:

\begin{itemize}
\tightlist
\item
  \textbf{Worldsheet CFT}: Virasoro emergence from modular data, modular invariance from overlap consistency
\item
  \textbf{Complete massless spectrum}: metric + B-field + dilaton + RR sectors, anomaly cancellation from worldsheet beta functions
\item
  \textbf{Internal sector matching}: the SM gauge content with \(c_{\text{int}} \approx 9\) is numerically consistent with integer WZW levels \(k = 7\) or \(8\) (§8.3)
\end{itemize}

The graviton is already contained implicitly via the Einstein equation derived from entanglement equilibrium.

\begin{center}\rule{0.5\linewidth}{0.5pt}\end{center}

\section{10. Summary}\label{10-summary}

OPH contains string theory as an emergent description of its edge/boundary degrees of freedom. The derivation chain is:

\begin{enumerate}
\def\labelenumi{\arabic{enumi}.}
\tightlist
\item
  \textbf{Edge-sector weights are exactly 2D Yang-Mills heat kernels} (Theorem 4.1)
\item
  \textbf{2D Yang-Mills admits a worldsheet expansion} via Gross-Taylor
\item
  \textbf{OPH supplies worldsheet kinematics}: conformal structure from Conf\((S^2) \cong\) SO\(^+(3,1)\), modular sewing from Markov collar structure
\item
  \textbf{OPH supplies target-space dynamics}: Einstein equation from entanglement equilibrium
\item
  \textbf{OPH supplies higher gauge structure}: 2-group gluing cocycles provide the B-field slot
\end{enumerate}

OPH\textquotesingle s edge-sector structure \textbf{is} a 2D gauge theory with exactly the heat-kernel structure that admits a worldsheet rewriting. This is not analogy , it is mathematical identity.

\begin{center}\rule{0.5\linewidth}{0.5pt}\end{center}

\section{Appendix A: Entanglement First Law Numerical Check}\label{appendix-a-entanglement-first-law-numerical-check}

The "entanglement first law" step (\(\delta S = \delta\langle K \rangle\)) is load-bearing in the Einstein derivation. Here we verify it numerically.

\subsection{A.1 Setup}\label{a1-setup}

Generate random bipartite density matrix on \(4 \times 4\) Hilbert space (2 qubits × 2 qubits). Perturb by small \(\varepsilon\) and compare:

\[\Delta S_A := S(\rho_A(\varepsilon)) - S(\rho_A(0))\] \[\Delta\langle K_A \rangle := \text{Tr}[(\rho_A(\varepsilon) - \rho_A(0)) K_A]\]

where \(K_A = -\log \rho_A(0)\).

\subsection{A.2 Expected Result}\label{a2-expected-result}

As \(\varepsilon \to 0\): \[\frac{|\Delta S_A - \Delta\langle K_A \rangle|}{|\Delta\langle K_A \rangle|} \to 0\]

linearly in \(\varepsilon\) (first-order identity).

\subsection{A.3 Code}\label{a3-code}

See \texttt{code/entanglement\_first\_law.py} for implementation.

\begin{center}\rule{0.5\linewidth}{0.5pt}\end{center}

\section{Appendix B: Key Equations Summary}\label{appendix-b-key-equations-summary}

{\def\LTcaptype{none} % do not increment counter
\begin{longtable}[]{@{}>{\RaggedRight\arraybackslash}p{0.480\textwidth}>{\RaggedRight\arraybackslash}p{0.480\textwidth}@{}}
\toprule\noalign{}
Result & Equation \\
\midrule\noalign{}
\endhead
\bottomrule\noalign{}
\endlastfoot
Generalized entropy & \(S_{\text{gen}} = \frac{A}{4G} + S_{\text{bulk}}\) \\
Einstein from equilibrium & \(G_{ab} + \Lambda g_{ab} = 8\pi G \langle T_{ab} \rangle\) \\
Edge-sector weights & \(p_R(t) \propto d_R \, e^{-t C_2(R)}\) \\
Edge partition function & \(Z_{\text{edge}}(t) = \sum_R d_R^2 \, e^{-t C_2(R)} = K_t(1)\) \\
2D YM partition function & \(Z_{\text{2D YM}}(S^2) = \sum_R d_R^2 \, e^{-\frac{g^2 A}{2} C_2(R)}\) \\
Parameter identification & \(t = \frac{g_{\text{YM}}^2 A}{2}\) \\
String tension & \(T_s = \frac{1}{2\pi\alpha'} \sim \frac{1}{2\pi a_{\text{cell}}}\) \\
Genus expansion & \(\log Z = \sum_{h \geq 0} N^{2-2h} F_h(\lambda)\) \\
\end{longtable}
}

\begin{center}\rule{0.5\linewidth}{0.5pt}\end{center}

\section{References}\label{references}

This derivation uses:

\begin{enumerate}
\def\labelenumi{\arabic{enumi}.}
\tightlist
\item
  OPH axiom set (A1-A4) and regulator premises (R0-R1) from Part I of this manuscript
\item
  Heat-kernel / edge-sector results from Part V of this manuscript
\item
  Gross-Taylor 2D YM / string duality (Gross \& Taylor, 1993)
\item
  Jacobson entanglement equilibrium (Jacobson, 1995, 2016)
\item
  Standard results: Tomita-Takesaki modular theory, Peter-Weyl theorem, Raychaudhuri equation
\end{enumerate}
